\section{Private equity and private debt}

\begin{qfdefinition}[Private equity: features, stages and strategies, how funds operate, and its investment characteristics]\label{def:private-equity-features}
\qffield{Statement} \emph{Features.} Private equity buys stakes in companies that are not listed, or delists them, through funds that are limited partnerships with a life of about ten years \citep[pp.~594--595]{s7-ang}. Four characteristics distinguish it \citep[pp.~5--7]{metrick-yasuda-survey-2011}:
\begin{enumerate}
\item the managers are financial intermediaries investing other people's money;
\item they invest in private companies, about which public information is scarce, so holdings are not marked to market;
\item they monitor and help manage the companies, through board seats, veto rights and contingent control rights;
\item they must exit, by an IPO, a sale to another investor or a sale to a company, to return cash within the fund's life.
\end{enumerate}

\emph{Strategies} (Ang, pp.~595--596): \emph{venture capital} takes minority stakes in young companies with little or no leverage; \emph{buyouts} acquire control of mature companies with high leverage, and represent most of the money; \emph{growth} capital takes minority stakes in expanding profitable companies; distressed and turnaround funds buy troubled companies; secondary funds buy existing fund stakes; LPs also co-invest.

\emph{Value creation in buyouts.} Kaplan and Str\"omberg group the changes private-equity firms make as \emph{financial engineering} (management equity and leverage, which disciplines the use of free cash flow), \emph{governance engineering} (active, smaller boards) and \emph{operational engineering} (industry and operating expertise) \citep[NBER WP 14207, pp.~13--16]{pe-kaplan-stromberg-2009}. Surveyed buyout managers rank growth in the business and operational improvements above leverage and multiple arbitrage as return drivers \citep[NBER WP 21133, p.~27]{w5-gompers-kaplan-mukharlyamov}.

\emph{Investment characteristics.}
\begin{itemize}
\item \emph{Illiquidity and the J-curve}: capital is called over years and returned on exits; early fees and investments at cost make interim returns negative \citep[pp.~16--18]{ippolito-private-capital-2020}.
\item \emph{Returns and dispersion}: in Burgiss data on nearly 1,400 US funds, average buyout funds beat the S\&P~500 by 20--27 percent over a fund's life, more than 3 percent a year; venture funds outperformed in the 1990s and underperformed in the 2000s \citep[abstract]{pe-harris-jenkinson-kaplan-2014}. Returns vary widely across funds and persist from one fund of a manager to the next \citep[NBER WP 9807, pp.~1--2]{pe-kaplan-schoar-2005}.
\item \emph{Cyclicality}: commitments and buyout activity rise in booms, when expected returns are lowest \citep[p.~596]{s7-ang}; \citep[p.~3]{pe-kaplan-stromberg-2009}.
\item \emph{Holding period}: the median buyout is held about six years; 38 percent of exits are sales to strategic buyers, 24 percent secondary buyouts, 14 percent IPOs (pp.~11--12).
\end{itemize}

\qffield{Practice} \emph{Ashby Fund I} (original, \hyperref[meth:private-fund-multiples-pme]{metrics}): four buyouts bought in years 0--3 are worth 0.93 of paid-in capital after two years, because of fees and holding at cost. Exits in years 4--7 lift the investor's multiple to 1.52, with the since-inception IRR rising from $-14.8$ percent to 11.7 percent.

\qffield{Conditions} The return evidence covers US funds over specific vintages and is net of fees but not adjusted for leverage or market beta beyond the S\&P~500 comparison. The NBER working-paper page numbers are used for the papers marked so.

\qffield{Pitfalls} Treating average performance as attainable when dispersion and access dominate outcomes. Reading interim IRRs early in a fund's life. Ignoring that leverage raises the risk of buyout equity.

\qffield{See also} \hyperref[meth:private-equity-strategies]{strategies over the life cycle}; \hyperref[def:private-debt-features]{private debt}; \hyperref[meth:private-capital-diversification]{diversification}.

\end{qfdefinition}

\begin{qfdefinition}[Private debt: direct and unitranche lending, mezzanine and distressed debt, their instruments and investment characteristics]\label{def:private-debt-features}
\qffield{Statement} \emph{Features.} A private debt fund lends to, or buys the debt of, private companies, as a limited partnership whose commitments are called over the investment period \citep[p.~13]{ippolito-private-capital-2020}. Private debt funds lend when banks do not, and when the risk does not suit private equity \citep[p.~127]{demaria-private-markets-2021}. The main strategies (pp.~127--136):
\begin{itemize}
\item \emph{Direct and senior lending}: senior loans to small and mid-sized companies, often for one-off purposes such as acquisitions, analysed by due diligence rather than bank scoring. Bank retreat after the 2008 crisis and Basel~III expanded it. Gross expected IRRs are 5--8 percent in Western Europe (4--6 percent net) and 7--9 percent unlevered in the US (pp.~133--134).
\item \emph{Unitranche}: one facility combining senior and subordinated debt at a blended rate, with ranking among lenders set by an agreement among lenders; often bullet, five to seven years, up to about 5.5 times EBITDA \citep[pp.~43--44]{ippolito-private-capital-2020}.
\item \emph{Mezzanine}: subordinated debt with cash and payment-in-kind interest and often warrants, contractually subordinated by an intercreditor agreement \citep[pp.~39--42]{ippolito-private-capital-2020}.
\item \emph{Distressed debt}: debt bought at a discount for trading, or to take control of the business through a restructuring (``loan-to-own'') \citep[pp.~127--128]{demaria-private-markets-2021}.
\end{itemize}

\emph{Investment characteristics.}
\begin{itemize}
\item Contractual, mostly floating-rate income with covenants and security, and target returns ranked by seniority: about 3--5 percent for senior secured debt, 6--8 percent for senior unsecured, 12--20 percent for mezzanine with warrants \citep[p.~21, Table~1.3]{ippolito-private-capital-2020}.
\item A shorter investment period (about three years) and life (about eight years) than equity funds, with early cash income and so a shallow J-curve \citep[pp.~61--62]{demaria-private-markets-2021}.
\item Risk lies in default and recovery rather than in growth; mezzanine funds rarely lose capital in aggregate, but individual vintages can (p.~133).
\item Illiquidity and limited price data, as for private equity.
\end{itemize}

\qffield{Practice} \emph{Corran Industrial} (original, \hyperref[meth:private-debt-credit-analysis]{credit analysis}): a buyout financed with 250 of debt at 5.0 times EBITDA can use a 9 percent unitranche or 175 of senior debt at 7 percent with 75 of mezzanine at 7 percent cash plus 5 percent PIK. The unitranche costs 0.5 points more but is one document with one lender group.

\qffield{Conditions} Return expectations are those reported by the sources for their periods and markets; leverage at the fund level raises them.

\qffield{Pitfalls} Treating private debt's smooth reported returns as low risk. Comparing a senior direct-lending yield with a mezzanine return without the difference in ranking.

\qffield{See also} \hyperref[meth:private-debt-financing-structures]{financing structures}; \hyperref[meth:private-debt-credit-analysis]{credit analysis}; \hyperref[meth:special-situations-distressed]{special situations}.

\end{qfdefinition}

\begin{qfmethod}[Diversification benefits of private capital: low reported correlations, smoothing and unsmoothing, cycle-dependent strategies, and the opposing views]\label{meth:private-capital-diversification}
\qffield{Statement} Private equity and debt add exposures that public markets lack: companies and loans that are not listed, control-oriented strategies, and strategies that do well at different points of the cycle. Demaria and co-authors argue that venture capital does best when invested at the end of a recession, buyouts when assets are bought just after a peak, and distressed debt when equity prices and growth fall; distressed debt is ``anti-correlated'' with venture capital \citep[pp.~162--163]{demaria-private-markets-2021}.

\emph{Measured correlation.} Reported returns are appraisals, which lag and smooth true values. With an AR(1) smoother $r^*_t=(1-\phi)r_t+\phi r^*_{t-1}$, the reported variance is $(1-\phi)/(1+\phi)$ of the true variance, and contemporaneous betas and correlations shrink in proportion (Derivation). Ang concludes that reported illiquid-asset returns ``are not returns''; volatilities, correlations and betas from them are too low. The Geltner and Ross--Zisler correction inverts the smoother \citep[pp.~416--421]{s7-ang}:
\[ r_t=\frac{r^*_t-\phi\,r^*_{t-1}}{1-\phi}. \]
Getmansky, Lo and Makarov model hedge-fund serial correlation the same way, as a sign of illiquidity exposure \citep[NBER WP 9571, abstract]{wide-getmansky-lo-makarov}.

\emph{The opposing view.} Demaria and co-authors call de-smoothing buyout returns a ``myth''. Appraisals of control stakes reflect rights, warranties and premiums that no listed price captures, so more frequent proxies add noise, and fund-of-funds return distributions show low loss rates \citep[pp.~160--166]{demaria-private-markets-2021}. The evidence on cash-flow performance shows buyouts beating public markets on average net of fees \citep[abstract]{pe-harris-jenkinson-kaplan-2014}, but not that their risk is low.

\qffield{Practice} \emph{Smoothing simulated} (original, 40,000 quarters, seed fixed). True quarterly returns with volatility 10 percent and market beta 1.0 are reported through a smoother with $\phi=0.6$:
\begin{itemize}
\item reported volatility is 5 percent and contemporaneous beta 0.40, so the reported Sharpe ratio doubles;
\item the estimated autocorrelation is 0.60, and unsmoothing restores volatility and beta to within 0.1 percent.
\end{itemize}
An allocator who uses the reported figures would see a diversifier with half its true risk. The cash-flow PME of a fund is unaffected by smoothing, because it uses only calls and distributions (\hyperref[meth:private-fund-multiples-pme]{multiples and PME}).

\qffield{Conditions} The AR(1) smoother is a model; real appraisal lags vary by asset and manager. The cycle claims are the authors' interpretation of patterns, not tests.

\qffield{Pitfalls} Optimising an allocation on appraisal-based volatilities and correlations. Concluding that unsmoothed returns are the truth rather than another estimate. Treating cycle diversification as available to an investor who commits pro-cyclically.

\qffield{See also} \hyperref[meth:private-public-market-comparison]{private against public markets}; \hyperref[meth:alternatives-allocation-approaches]{allocation}.

\end{qfmethod}
